% ============================================================
%  ABSTRAK (Bahasa Indonesia & Inggris)
%  Pedoman: TNR 12pt, single spacing, maks 200 kata per bahasa
% ============================================================

\phantomsection
\thispagestyle{chapterpage}
\addcontentsline{toc}{chapter}{ABSTRAK}

{\singlespacing

%% ---- Abstrak Bahasa Indonesia ----
%\begin{center}
%  {\bfseries ABSTRAK}
%\end{center}
%
%\noindent
%\NamaMahasiswa, \NPM \\
%\MakeUppercase{\JudulPISatu}\\
%Penulisan Ilmiah. Informatika. Fakultas Teknologi Industri. Universitas Gunadarma. \Tahun
%
%\noindent
%Kata Kunci: \KataKunci\\
%(xii + 62 + Lampiran)
%
%\medskip
%\noindent
%% Isi abstrak – maksimum 200 kata, 1 paragraf
%Color Vision Deficiency (CVD), commonly known as color blindness,is a genetic condition that affects a person's capacity ability to differentiate between certain light wavelengths.
%Although eyesight itself is usually normal, individuals with CVD find it difficult to distinguish colors within the red, green, and blue spectrums.
%The most common type is Dichromacy, which includes Protanopia (red-blind), Deuteranopia (green-blind), and Tritanopia (blue-blind).
%Existing assistive technologies, such as color-filtering glasses, are designed to separate overlapping red and green wavelengths, help with color distinction but are often expensive.
%Additionally, many existing smartphone applications require users to upload images for processing, which prevents real-time operation and limits immediate visual feedback.
%This paper proposes the implementation of the LMS Daltonization algorithm on the Android platform using real-time image processing.
%The method converts the camera's input from the RGB color space to the LMS color space.
%This allows for the mathematical simulation and adjustment of the missing color spectrums.
%The application is expected to help users distinguish objects with specific color spectrum through a smartphone, with low latency that supports real-time usage.
%
%\noindent
%Daftar Pustaka (<<tahun awal>> -- <<tahun akhir>>)
%
%\vspace{2\baselineskip}

% ---- Abstract Bahasa Inggris ----
\begin{center}
  {\bfseries ABSTRACT}
\end{center}

\noindent
\NamaMahasiswa, \NPM \\
\MakeUppercase{\JudulPISatu}\\
Scientific Writing. Informatics. Faculty of Industrial Technology. Gunadarma University. \Tahun

\smallskip
\noindent
Keywords: \Keywords\\
(xii + 62 + Appendix)

\medskip
\noindent
% Terjemahan abstrak ke Bahasa Inggris – maksimum 200 kata
Color Vision Deficiency (CVD), commonly known as color blindness,is a genetic condition that affects a person's capacity ability to differentiate between certain light wavelengths.
Although eyesight itself is usually normal, individuals with CVD find it difficult to distinguish colors within the red, green, and blue spectrums.
The most common type is Dichromacy, which includes Protanopia (red-blind), Deuteranopia (green-blind), and Tritanopia (blue-blind).
Existing assistive technologies, such as color-filtering glasses, are designed to separate overlapping red and green wavelengths, help with color distinction but are often expensive.
Additionally, many existing smartphone applications require users to upload images for processing, which prevents real-time operation and limits immediate visual feedback.
This paper proposes the implementation of the LMS Daltonization algorithm on the Android platform using real-time image processing.
The method converts the camera's input from the RGB color space to the LMS color space.
This allows for the mathematical simulation and adjustment of the missing color spectrums.
The application is expected to help users distinguish objects with specific color spectrum through a smartphone, with low latency that supports real-time usage.

\smallskip
\noindent
References (<<start year>> -- <<end year>>)

} % end singlespacing

\newpage
